\section{The endpoint in rough domains}
\label{sec:q4rough}

We work at $q=4$ throughout this section.  We use the whole-plane
height convergence of Proposition~\ref{q4field:wholeplane-input}
and its regular-domain form, Theorem~\ref{q4field:regular-domain}.
Write
\[
 \sigma^2=\frac{2}{\pi},\qquad \lambda_0=\sqrt{\frac{\pi}{8}},
 \qquad \sigma\lambda_0=\frac12.
\]
The normalization of the Dirichlet Gaussian free field is
$G_D=(-\Delta_D)^{-1}$, so its Green function has singularity
$(2\pi)^{-1}\log(1/|z-w|)$.  A signed height has an independent fair
unit jump across each closed contour.  Its offset is specified separately.
Our purpose is to identify the conditional field in domains cut out by
the contours themselves.  Their boundaries need not be Jordan curves at
this stage.

\subsection{Pockets and their conditional laws}

A \emph{wall of color $C$} means an actual connected open wall of that
color, rather than an identification of separated vertices.  Denote the
other color by $C'$.  All the following discrete drawings may be separated
within a distance $O(\delta)$, where $\delta$ is the mesh, so that their
contours are simple and have the prescribed traversal order.

\begin{lemma}[Exterior-readable pockets]
\label{q4pocket:discrete-pockets}
In a homogeneous $C$-wall disk, each outer contour is the external
perimeter of an outermost actual $C'$-cluster.  Behind it lie pockets
bounded by the maximal simple $C'$-open cycles, ordered by inclusion of
their interiors.  Their interiors are disjoint; several pockets can
belong to one outer contour.  The contours, these cycles, their
assignment to contours, and all switches outside the pocket interiors
can be determined without querying strictly interior switches.
Conditional on these data, the pocket interiors have independent
homogeneous $C'$-wall laws.

In a two-arc disk, first reveal the distinguished path.  On either side,
its bank together with that side's wall arc determines analogous
maximal-cycle pockets of the wall color.  Conditional on the path and
the exterior pocket data, all these pockets have independent
homogeneous laws.  A point wound by a containing contour or side
closure, and farther than $O(\delta)$ from its bank walk, lies in one of
the corresponding pockets.  Every further contour winding that point
is contained in that pocket.
\end{lemma}

\begin{proof}
Draw each cluster as a small regular neighborhood of its actual graph.
Its boundary follows the paired medial pieces within $O(\delta)$.
Every perimeter behind its external perimeter lies in a bounded face
of that graph.  Such a face is enclosed by a simple cycle.  Equivalently,
the bank of the external perimeter is a closed graph walk; decomposing
this walk into simple cycles shows that every wound point with the
specified clearance belongs to a maximal cycle interior.

Two maximal cycle interiors cannot overlap without one containing the
other.  Otherwise their boundaries meet at two or more vertices, their
union is contained in one biconnected block, and the simple outer boundary
of that block encloses both.  This contradicts maximality.  Shared
vertices or edges cause no difficulty, since the open interiors remain
disjoint.

We justify the assertion about revealing information.  Fix a possible
list of cycles and expose its wall edges and everything outside its
open disks.  Any connection between vertices of one cycle that uses
its interior already has a connection along that cycle.  More strongly,
adding arbitrary interior edges cannot join two distinct exterior
biconnected blocks.  After removal of an exterior cut vertex, the
attachments belonging to any one wall cycle remain connected along
the rest of that cycle.  Thus an interior addition cannot repair an
exterior cut vertex.  The exterior block containing a maximal cycle
already has all its vertices in that cycle's closed disk: otherwise
the simple external boundary of the block would be a larger cycle.
A new cycle that both enters and leaves a listed disk would have two
distinct wall attachments and would consequently belong to its wall
block.  It cannot enlarge that block outside the disk.  The same
argument applies if several listed disks are traversed, since they
would then have to belong to the same exterior block.  Nor can internal
changes alter an external cluster perimeter.  Consequently the event
specifying the proposed list and its exterior is measurable from the
revealed edges.  Each wall is actually connected, so the domain Markov
property gives the asserted product of homogeneous laws.

For the marked path, close each side using its wall arc and the
appropriate half-sides of the terminal gates.  The same-color bank and
wall form a closed walk of actual open edges.  The path lies outside
every cycle made from this walk, so its exploration queries no switch
strictly inside such a cycle.  The preceding argument therefore applies
on both sides.  Finally, a further contour leaving a pocket and winding
a point inside it would have to surround its wall cycle.  This is
impossible because that cycle belongs to an actual cluster connected
to the original wall.
\end{proof}

For a family $\mathcal D$ of disjoint homogeneous disks, let
$H_{\mathcal D}$ be the signed height with zero offset and independent
laws in the disks.  Its value at a point refers to the disk containing
that point.

\begin{lemma}[Comparison of zero-offset moments]
\label{q4pocket:moment-comparison}
Let $\mathcal D_-$ and $\mathcal D_+$ be two families of homogeneous
discrete disks of the same wall color.  Suppose that each disk of
$\mathcal D_-$ is contained in a disk of $\mathcal D_+$.  Assume all
disks lie in a fixed bounded region.  For distinct points
$z_1,\ldots,z_m$ in the smaller interiors,
\begin{equation}
 \mathbb E\prod_{j=1}^m H_{\mathcal D_-}(z_j)
 \ \leq\
 \mathbb E\prod_{j=1}^m H_{\mathcal D_+}(z_j)+o(1).
 \label{q4pocket:moment-order}
\end{equation}
The error tends to zero with the mesh, uniformly when the points have
fixed positive pairwise separation.  The statement compares moments, not signed
heights in an ordered coupling.
\end{lemma}

\begin{proof}
On graphs of the opposite color, the smaller laws are independent free
random-cluster laws.  Their product is stochastically dominated by the
restriction of the larger laws: reveal the remaining edges and use
the induced boundary partitions, which dominate free boundary
conditions.  Couple these opposite-color configurations by inclusion.

Consider a smaller outer contour that is relevant to the test points.
Its cluster is included in a larger-model cluster, but the latter may
itself be surrounded by a different cluster.  Select the \emph{outermost
larger contour containing the filled hull} of the smaller cluster.
One must make this selection even when the two clusters do not merge.
Using nested regular neighborhoods of included graphs, or the
cycle decomposition in Lemma~\ref{q4pocket:discrete-pockets}, shows
that this selected contour contains all the wound test points with
clearance.  Every relevant smaller maximal cycle lies in a maximal
cycle behind the selected larger contour.  In the surrounding case,
the smaller cycle lies in a bounded face of the outer cluster and is
therefore enclosed by one of its maximal cycles.

Retain only the two sets of exterior and pocket data from this
coupling.  Now resample their interiors using the conditional laws
of Lemma~\ref{q4pocket:discrete-pockets}.  Their wall color has reversed,
so the same free opposite-color comparison applies at the next stage.
Iterate.  This is a sequence of couplings of conditional marginals;
it does not claim monotonicity of nesting depth in the original
edgewise coupling.  At every fixed number of stages it gives inclusion
at the same numerical nesting levels.  A smaller contour's group of
test points can merge with another such group, but cannot split.
There is no extra contour of that level between a selected outer
contour and a pocket behind it: the pocket wall is part of that very
contour's cluster.

To translate this construction into moments, expand each height as
the sum of signs of its surrounding contours and average the signs.
A term survives exactly when every assigned contour receives an even
number of the indices $1,\ldots,m$.  Encode that term by the vector of
nesting levels $(\ell_1,\ldots,\ell_m)$ assigned at the respective
points.  The same vector is admissible in the larger model: mergers
add even multiplicities and hence preserve their parity.  Distinct
vectors remain distinct assignments.  This is an injection between
the terms counted by the two sign-averaged moments.

For a surviving term, each assigned contour winds at least two distinct
test points and has diameter bounded below by their separation.  All
its ancestors have that lower bound as well.  Truncate the number of
such contours and ancestors.  The free-bank arm bound makes the event
that any of these contours passes within $O(\delta)$ of a test point
negligible.  On its complement all the required pocket containments
hold.  Lemma~\ref{q4field:free-bank} and the higher count moments
in Lemma~\ref{q4field:moment-kernels} remove the exceptional event
and then the truncation by uniform integrability.  This proves
\eqref{q4pocket:moment-order}.
\end{proof}

\subsection{Geometry before identification}

The arm and passage estimates give three properties of every
subsequential contour limit that we shall use.  We spell out their
endpoint content to avoid importing a crossing estimate at a true
$q=4$ boundary.

\begin{lemma}[Status disks and boundary traversal]
\label{q4pocket:geometry}
For every positive-diameter limiting contour, and for the limiting
distinguished path, the set of distinct multiple visits in the
original-domain interior has Hausdorff dimension at most zero.
Deterministic interior points are almost surely avoided.  At every
interior trace point there are disks of each winding status arbitrarily
nearby.  In the marked case the statuses are read after closing the
path on either wall arc.  No nonconstant subtraversal of a limit lies
entirely on the original Jordan boundary.

These statements and curve tightness hold both for the regular wall
experiments and for the marked Jordan approximations in the theorem.
\end{lemma}

\begin{proof}
The first two assertions follow from
Proposition~\ref{q4field:four-arm} and
Corollary~\ref{q4field:contact-dimension}, together with the one-arm
estimate.  For the two-status assertion use the attached-circuit
argument of Section~\ref{sec:loops}, with both banks retained.  Here
is the only modification required at $q=4$.  Enlarge the first
crossing band a little laterally, and give a label only to an actual
component containing a crossing of the original core.  After a
component has been searched and removed, comparison with full wiring
on the enlarged band's boundary is still valid.  A buffered
opposite-color transversal extending past both lateral ends of the
core blocks any further core crossing.  Bulk RSW gives this blocker
a uniformly positive probability.  The number of labels is therefore
tight, and each label represents an actual connection in the revealed
past.  The leftmost crossings from a label's bottom sources in
successive core slabs can now be used exactly as in the attached-circuit
argument; the fresh-side circuits lie in interior balls.  For a
distinguished path, close it along the remote wall when reading the
two statuses.  The argument also works with exposed bulk data outside
the test patch.  Only buffered bulk RSW is used
\citep{DCLM2019Isoradial}.

For the boundary assertion, fix a deterministic boundary point away
from the marks.  A macroscopic passage approaching it supplies an
actual arm of the bank color that is not wired there.  In local
annuli this arm is bounded by the free-bank comparison.  Along a
primal wired arc the dual bank uses internal square faces and has no
dual exit through that arc.  Along the complementary arc the primal
bank uses actual sampled edges and has no primal ties.  Exterior
dual identifications occur only beside the complementary arc.
Turns can be followed by adjacent actual bank edges within
$O(\delta)$.  Any other induced ties in a buffered annulus away from
the other arc attach the surviving graph only at its rims.  Thus the
same arm bound applies in the specified Jordan approximations.

A nonconstant boundary subtraversal would cover an arc and hit a
point of a fixed countable dense subset away from the marks, by a
macroscopic passage.  The preceding bound excludes this.  For
tightness, at each displacement scale use a finite cover of
subpassage tests away from the two marks.  Every displacement contains
a smaller displacement a fixed fraction of that scale away from
both marks.  Marked uniform Jordan convergence and separation of
the two boundary arcs allow its annulus to be cropped in a
neighborhood missing at least one arc.  The same passage estimates
apply there.  Interior no-spur and multiple-visit tests always retain
positive clearance from the original boundary.
\end{proof}

\subsection{Protected cuts and pointed kernels}

We next give the deterministic approximation needed for the conditional
field.  A single limit contour may have several open winding components;
filling in its hull would incorrectly join some of them.

\begin{lemma}[Outer comparisons with protected cuts]
\label{q4pocket:protected-cuts}
Let $L_n$ be simple discrete contour closures in a bounded Jordan
domain $V$, converging uniformly modulo reparametrization to $L$.
The closures may include one original wall arc.  Assume the geometric
properties of Lemma~\ref{q4pocket:geometry}.  Let $T$ be the trace
of $L$, and let $U$ and $E$ be, respectively, the open sets of winding
status one and zero off $T$, with infinity in $E$.  Then the
components of $U$ are simply connected pointed kernels of the
interiors of $L_n$.

For every $\beta>0$ there exist simple regular comparison polygons
$B_j$, and protecting balls centered on $T\cap V$ or $\partial V$,
such that:
\begin{enumerate}[label=\textup{(\roman*)}]
\item the pointed kernels of $B_j$ at points of $U$ are precisely
the separate components of $U$;
\item every compact subset of $U$ is eventually contained in $B_j$;
\item at each fixed $j$, deleting the part of a sufficiently fine
discrete pocket outside $B_j$ changes edges only in finitely many
protecting balls, of radii $r_i$, with
$\sum_i r_i^\beta\longrightarrow0$ as $j\to\infty$;
\item $B_j$ is connected to a region outside $\overline V$.
\end{enumerate}
The finite protecting selection is made after constructing the cuts.
\end{lemma}

\begin{proof}
Every point of $T\cap V$ is approached by $E$, by the two-status
property.  Points of $T\cap\partial V$ are approached by the exterior
of $\overline V$.  Hence
\begin{equation}
 \overline E=\mathbb C\setminus U.
 \label{q4pocket:exterior-closure}
\end{equation}
Compact subsets of $U$ persist in the approximating interiors by
winding stability.  Conversely, every excluded point is arbitrarily
close to persistent exterior points, by
\eqref{q4pocket:exterior-closure}.  Thus the component containing a
fixed point of $U$ is exactly its pointed kernel.  Kernels of simply
connected domains are simply connected.

Let $M\subset T\cap V$ be the distinct multiple-visit set, after
constant pauses have been erased.  Given $\eta>0$, cover \emph{all}
of $M$ by a countable collection of open balls $D_i$, centered on
$M$, with
\[
 \sum_i\operatorname{rad}(D_i)^\beta<\eta,
 \qquad \sup_i\operatorname{rad}(D_i)<\eta.
\]
This is possible because $\dim_{\mathrm H}M\leq0$.  Set
$G=E\cup\bigcup_iD_i$.

We first show that any $x\in E\cap V$ can be joined to the exterior
through $G$, after adding at most one arbitrarily small ball centered
on $\partial V$.  If the closure of the path component of $x$ in
$G$ reaches $\partial V$, such a ball connects it to the Jordan
exterior.  Otherwise that component has a compact frontier
$S\subset V\setminus G$ separating $x$ from infinity.  The set
$S$ meets $T$ only at unique-preimage points; its other points lie
in $U$.

There are polygonal closed walks, arbitrarily close to $S$, of
nonzero winding about $x$.  To see this, grid the bounded component
of $\mathbb C\setminus S$ containing $x$, retain its edge-adjacent
squares reachable from $x$ without meeting $S$, and trace their
external perimeter.  We can lift these walks, with an error tending
uniformly to zero, into the closed disks bounded by $L_n$.  Near a
point of $S\cap U$, a fixed interior ball persists.  Near a point
$z\in S\cap T$, all boundary parameters approaching $z$ tend to
its unique preimage under $L$.  Nearby points of the approximating
closed disk can therefore be connected near $z$: first go by a
closest segment to the boundary, and then use the short boundary
arc between the resulting nearby parameters.  The closest segments
stay in the closed disk.  The errors are uniform along the walk;
otherwise compactness would give one of the two local situations
just considered and contradict failure of uniformity.  The lifted
closed walk still winds about $x$, whereas $x$ lies outside the
approximating disk.  This contradicts contractibility of that disk
and proves the escape assertion.

Choose a large square $Q$ containing $\overline V$ in its interior.
Take finite increasingly fine nets of $\overline E\cap Q$ whose
vertices belong to $E$.  Join each net vertex to an exterior collar
of $Q$ using the escape assertion.  The finitely many additional
boundary-centered balls can have arbitrarily small total
$\beta$-content.  Select a finite polygonal tree joining these
paths to one root on $\partial Q$, otherwise lying in the open
allowed region; gridding the paths and deleting cycles gives such
a tree with axis-parallel edges.  Remove a sufficiently thin
polygonal neighborhood of this tree, attached to $\partial Q$.
Using narrow rectangular links and square junctions leaves one
simple rectilinear polygon $B$.  Keep an open part of the collar
outside $\overline V$ in $B$.

As the nets become finer, the removed tree approaches every point
of $\overline E\cap Q$.  Any part of it entering $U$ lies in the
protecting balls, whose radii tend to zero and whose centers lie
on $T$ or $\partial V$.  Hence compact subsets of $U$ persist,
while the pointed kernels cannot extend through an excluded point.
This proves the first two assertions.

For the finite selection, fix the tree and its closed thin
neighborhood.  It is compact and lies in the union of $E$ and
the countable protecting balls.  A finite subcover uses compact
subsets of $E$ and finitely many of those balls, with a positive
margin.  The compact subsets of $E$ are disjoint from the closed
disks bounded by $L_n$ for all large $n$.  Thus only the finitely
many selected balls can contain cuts affecting the discrete
pockets.  Their total content is no larger than that of the full
cover.  This is why the countable cover is used before selecting
the tree.  A further $O(\delta)$ adjustment makes the comparison
polygon a regular wall polygon on the appropriate color lattice,
without losing the protecting margins.
\end{proof}

\subsection{Transfer of the conditional field}

\begin{proposition}[Rough-interior moment transfer]
\label{q4pocket:rough-transfer}
Condition on exterior pocket data converging to a containing contour
as in Lemma~\ref{q4pocket:protected-cuts}.  At every fixed distinct
tuple of points in $U$, the joint moments of the residual zero-offset
height converge to the Wick moments of independent fields
$\sigma\Phi^0_\Omega$ in the components $\Omega$ of $U$.
This holds along every such convergent sequence of exterior data.
Several containing contours at one generation may be treated
simultaneously, taking the product over them.
\end{proposition}

\begin{proof}
Write $\beta>0$ for the power in the free-bank arm bound.  Fix a
comparison polygon $B$ and the finite protecting balls supplied
by Lemma~\ref{q4pocket:protected-cuts}.  In every pocket delete
opposite-color edges outside the lattice approximation of $B$.
Equivalently, intersect its set of elementary wall-color cells
with the corresponding cell set of $B$, retaining only
opposite-color edges between surviving adjacent cell centers.
All changes are in the protecting balls.

The original opposite-color laws are product free laws.  Couple
them with the modified laws by exploring, in the larger law, the
actual clusters incident to the deleted edges.  Once those
clusters are exposed, all discrepancies are shielded and the
unqueried graphs have the same conditional law.  The argument
is the same as the gap-closing coupling in
Section~\ref{sec:q4regular}.  If the ball radii are $r_i$ and $e>0$
is fixed, the probability that a discrepancy reaches farther than
$e$ from their centers is at most
\begin{equation}
 C\sum_i(r_i/e)^\beta.
 \label{q4pocket:cut-error}
\end{equation}
Indeed every such discrepancy supplies an actual opposite-color
arm in the original product free law, dominated by the full-lattice
comparison.  Contours through changed switches have banks in
these discovered clusters and obey the same locality bound.
Choose $e$ smaller than the tuple's separation and its clearance
from $T\cup\partial V$.  Every sign-averaged residual moment term
uses a contour of at least the pair-separation diameter.
The macroscopic count moments and H\"older's inequality therefore
make the moment error tend to zero with the total protecting
content in \eqref{q4pocket:cut-error}.

We verify the laws and offsets of the retained regions.  Each
side-connected component of the retained wall-color cells is
again a simple polygon disk.  Every globally missing cell has
a path of missing cells to the exterior: it is outside at least
one of the two simple grid polygons being intersected.  Thus a
retained component has no hole.  Its boundary cannot pinch at
a vertex with two opposite retained cells in that component
and two opposite missing cells.  The retained pair would be
connected by a side-adjacent path, while the missing pair is
connected through the exterior; these two disjoint connections
between alternating sectors contradict planar separation.
Apply this argument separately to each old pocket.  The
retained free opposite-color graph hence factors into
independent homogeneous disk laws.

Each retained component also reaches an \emph{old} pocket wall.
Start from one of its cells and take a side-adjacent path in
$B$ to the part of $B$ outside $\overline V$.  The first exit
from the old pocket crosses an old wall edge; the path before
that exit remains in the same retained component.  The component
therefore has an actual wall-color connection to the old wall.
No exterior contour can surround it and add an unrecorded
height jump.  Its residual offset is zero.  Finally, a connected
compact neighborhood of any test point began in one old pocket
for large $n$: the bank walk tracking $L_n$ winds around that
compact, so Lemma~\ref{q4pocket:discrete-pockets} places it in
a pocket.

For the upper bound, apply
Lemma~\ref{q4pocket:moment-comparison} from the retained disks
into $B$, and then apply
Theorem~\ref{q4field:regular-domain}.  For the lower bound,
make no deletions.  Choose one regular polygon disk compactly
inside each involved component $\Omega$, containing its test
points, and compare their product laws into the original
pockets.  Such inner disks can exhaust each simply connected
$\Omega$ in the pointed-kernel sense.

Dirichlet Green functions converge along simply connected
pointed kernels.  Between points whose limiting kernels are
different, they tend to zero; this follows by conformal mapping
convergence, or by the vanishing probability of crossing the
separating bottleneck before exit.  Therefore the Wick kernels
of both the inner polygons and the protected outer polygons
converge to those of independent Dirichlet fields in the
components $\Omega$.  Take the mesh limit first at fixed
polygons and protectors; next improve the inner and outer
approximations and let their total protecting content vanish.
The two bounds coincide.  Already assigned ancestor signs, or
the marked path's side constants, multiply lower-order residual
moments and are restored afterwards.  They do not enter the
moment comparison.
\end{proof}

\section{Identification at the endpoint and the boundary convention}
\label{sec:q4identification}

\subsection{A global conditional decomposition}

Take a joint subsequential limit of the height and the distinguished
path in a regular two-arc polygon.  For a homogeneous initial square,
take a joint limit of the height and the full contour collection,
including nesting levels, parent relations, and independent fair signs.
These marked geometric limits can be obtained at countably many
positive-diameter cutoffs decreasing to zero, moving a cutoff when
necessary to avoid atoms.  At each cutoff there are only a tight
number of relevant contours and ancestors.  Thus finite marked
relations stabilize, and a diagonal extraction retains every positive
limiting contour.  The signs remain conditionally independent and
fair given the geometry, which was sampled without using them.

Fix a nesting level $k$ in the square.  Every deterministic interior
point lies almost surely in a positive winding region at that level.
Indeed alternating actual circuits on successively separated bulk
scales force arbitrarily many surrounding contours.  The domain
Markov bounds and bulk RSW make the failure probability tend to zero
as the number of scales grows, with the mesh limit taken first at
each finite selection of scales.  The same statement holds for
Lebesgue-almost every point by Fubini.  Winding regions of different
contours at a fixed level are disjoint, as their discrete predecessors
are.  Positive coincident traversal arcs would also contradict the
multiple-visit dimension bound.  In the marked case the two side
closures likewise cover almost every point away from their traces.

Let $\mathcal F_k$ denote the limiting exterior geometry and signs
through level $k$, or let $\mathcal F$ denote the limiting marked
path in the two-arc experiment.  Let $\mathcal O$ be the collection
of components of the corresponding open winding regions.  Define
$v$ on each component to be the sum of its containing signs through
level $k$, or its marked side value $0$ or $1$.  Its value on the
remaining area-zero set is immaterial.

\begin{proposition}[Conditional field in all complementary components]
\label{q4pocket:conditional-field}
Conditional on the stated limiting exterior data, the height has the
distributional decomposition
\begin{equation}
 h\ \stackrel{\mathrm{law}}=\
 v+\sigma\sum_{\Omega\in\mathcal O}\Phi^0_\Omega,
 \label{q4pocket:conditional-decomposition}
\end{equation}
where the Dirichlet fields are independent and extended by zero.
The identity holds as an identity of distributions on the whole
initial-domain interior, not merely after restriction to individual
components.
\end{proposition}

\begin{proof}
For a test function $f$, the independent sum on the right exists
in $L^2$, since Green-function monotonicity gives
\begin{equation}
 \sum_{\Omega\in\mathcal O}
 \iint_{\Omega^2}|f(z)|G_\Omega(z,w)|f(w)|\,dz\,dw
 \leq
 \iint_{V^2}|f(z)|G_V(z,w)|f(w)|\,dz\,dw.
 \label{q4pocket:covariance-domination}
\end{equation}
The same estimate, with a fixed negative Sobolev norm, defines the
sum as a random distribution.

Use a representation on which the extracted geometric data converge.
A fixed distinct tuple of deterministic interior points is almost
surely eventually covered, with positive limiting clearance, by the
relevant winding components.  Its offsets converge.  Conditional
on the discrete exterior pocket data, including the actual cycle
walls, its residual moments converge by
Proposition~\ref{q4pocket:rough-transfer}.  That proposition applies
along every realized convergent data sequence; it requires no
regularity of the unexplored cycle walls beyond the proved pocket
properties.

To pass to conditioning only on the stated limiting data, multiply
the discrete moment identities by a bounded continuous function of
those data.  The necessary uniform integrability follows from
Lemma~\ref{q4field:moment-kernels}.  At level $k$, terms already
assigned to ancestors have offsets bounded by $k$.  The remaining
sign-averaged terms require shared contours of at least the relevant
pair-separation diameter, and their counts have all polynomial
moments under the initial law.  The marked offset is bounded by
one.  Thus, for each $m$, a bound of the form
\[
 C_{m,k}\bigl(1+|\log d_*|\bigr)^m,
 \qquad
 d_* = \min\bigl\{1,|z_i-z_j|:i\ne j\bigr\},
\]
and its higher-moment versions apply also before coverage by the
limiting components has stabilized.  Such bounds are integrable
on tuples in compact subsets of $V$, including near diagonals.

Smear the identities against tensor products of smooth test
densities.  Fubini and uniform integrability pass every joint
moment to the limit, with every bounded continuous geometric
multiplier retained.  The almost-everywhere coverage removes the
remaining set from these integrals.  Higher even moments give
uniform integrability of the smeared variables themselves.
A countable determining family of test densities then yields the
conditional moment identities for distributions.  For each fixed
conditional geometry, the Gaussian law with the bounded offset
on the right of \eqref{q4pocket:conditional-decomposition} is
moment determining; equivalently, apply the one-dimensional
Gaussian moment criterion to every rational linear combination
of finitely many tests.  Hence the moment identities identify
the conditional law.  Because the tests range over the entire
initial domain, this argument also excludes a residual
distribution supported on the area-zero complement of the
components.
\end{proof}

\subsection{Local-set uniqueness with the exact height normalization}

We use two continuum identifications.  For the zero-boundary GFF,
the two-valued bounded-type thin local set with values $-a,b$ is
unique and field measurable when $a,b>0$ and
$a+b\geq2\lambda_0$.  At $a=b=2\lambda_0$ it is the non-nested
$\mathrm{CLE}_4$ gasket, with independent fair labels on its
holes \citep[Propositions~1 and~2 and Sections~4 and~6.3]
{AruSepulvedaWerner2019}.

For boundary data we use Theorem~3.21 of
\citet{AruLupuSepulveda2020}, with its immediately preceding
boundary criterion.  In the specialization needed here, the
initial domain is a disk, $u$ is bounded harmonic with finitely
piecewise constant boundary data in $[-a,b]$, and
$a+b\geq2\lambda_0$.  Uniqueness holds for bounded-type thin
local sets whose total conditional harmonic function is
constantly $-a$ or $b$ on each connected boundary component
of each complementary domain.  The interval here is closed:
boundary values equal to $-a$ or $b$ are permitted.  The
standard completed chordal $\mathrm{SLE}_4$ level line, for
boundary data $-\lambda_0,\lambda_0$ on the two arcs, has
exactly this conditional field decomposition
\citep{SchrammSheffield2013}.

\begin{lemma}[The limiting sets are bounded-type thin local sets]
\label{q4pocket:local-set}
For the first contour level in the square, let
\[
 A=\overline V\setminus\bigcup_{\Omega\in\mathcal O}\Omega.
\]
In the marked experiment let $A$ be the path trace together with
$\partial V$.  After the normalizations below, these are
bounded-type thin local sets.  In the homogeneous case their
harmonic values are $\pm2\lambda_0$.  In the marked case their
\emph{total} harmonic values are $\pm\lambda_0$.
\end{lemma}

\begin{proof}
In the homogeneous case, write $\Phi=h/\sigma$ and $m=v/\sigma$.
Theorem~\ref{q4field:regular-domain} says that $\Phi$ is a
zero-boundary GFF.  At the first level $v\in\{-1,1\}$, so
$m\in\{-2\lambda_0,2\lambda_0\}$.  In the marked case put
\[
 X=\frac{h-1/2}{\sigma}=\Phi+u,
 \qquad
 m=\frac{v-1/2}{\sigma}-u,
\]
where $\Phi$ is a zero-boundary GFF and $u$ is the bounded
harmonic function with boundary data $-\lambda_0,\lambda_0$
on the respective arcs.  Here $|m|\leq2\lambda_0$ and
$m+u\in\{-\lambda_0,\lambda_0\}$ on each component.

In either case Proposition~\ref{q4pocket:conditional-field}
gives the full conditional Dirichlet law of $\Phi-m$ in
$V\setminus A$.  This law depends on the exterior data only
through $A$.  Conditioning down to $(A,m)$ proves locality.
Moreover $m$ is an ordinary bounded function, integrated over
$V\setminus A$, and the distributional identity has no
additional part supported on $A$.  This is thinness.  The
uniqueness of a local set's harmonic distribution from the
field and set is also part of the local-set formalism
\citep[Lemma~2.3(1)]{AruLupuSepulveda2020}; no extra choice
of component labels is being added at this step.

For completeness, $A$ is connected and contains the entire
Jordan boundary.  Its complementary components are simply
connected by Lemma~\ref{q4pocket:protected-cuts}.  If $A$
were disconnected, split it into two nonempty disjoint compact
sets, putting the connected $\partial V$ in the first.  A fine
polygonal boundary around part of the second, at positive
distance from the first and from $\partial V$, would be a
closed walk in a component of $V\setminus A$ with nonzero
winding about a point of $A$.  This contradicts that
component's simple connectivity.  Thus $A$ is a nontrivial
connected compact set.  It has no isolated points, and no
component disjoint from $\partial V$.  These facts verify
the topological condition in the definition of a bounded-type
thin local set.  Conformal transport to a disk preserves all
these assertions.
\end{proof}

\begin{proposition}[Regular marked endpoint]
\label{q4pocket:regular-chordal}
The distinguished curves in the regular two-arc experiments
converge as ordered curves, up to increasing changes of parametrization,
to chordal $\mathrm{SLE}_4$ between the prescribed marks.
\end{proposition}

\begin{proof}
Apply the boundary-data uniqueness theorem with
$a=b=\lambda_0$.  The initial domain, after conformal transport,
has one complementary component, which is not a singleton.
The boundary function has two changes and takes values in the
closed allowed interval.  Its interior range meets
$(-\lambda_0,\lambda_0)$ by the strong maximum principle.
The height-gap condition holds with equality.  In every
component of $V\setminus A$, the total conditional mean is
constant $-\lambda_0$ or $\lambda_0$, so the stated boundary
criterion holds on every boundary component.  Finally
Lemma~\ref{q4pocket:local-set} supplies thinness, boundedness,
and the required topology.  All hypotheses of
\citet[Theorem~3.21]{AruLupuSepulveda2020} are therefore met.

The completed Schramm--Sheffield level line coupled to the
same field satisfies these same conditions.  Uniqueness
identifies the limiting closed set, together with its harmonic
labels, with that level line plus the original boundary.
In particular its interior trace is the simple chordal
$\mathrm{SLE}_4$ arc, which has no hits on the interiors of
the two boundary arcs.

This closed-set statement identifies the original ordered
curve as well.  Any extra boundary point off the identified
compact arc has a neighborhood disjoint from that arc; reaching
it would force a nonconstant boundary subtraversal, excluded
by Lemma~\ref{q4pocket:geometry}.  On a simple arc, any
genuine backtracking repeats a nondegenerate subarc.  Such a
subarc has positive Hausdorff dimension, contradicting the
dimension-zero multiple-visit bound.  After suppressing constant
pauses, the curve therefore traverses the arc once, in the
prescribed endpoint order.
\end{proof}

\subsection{Contours, labels, and all nesting levels}

\begin{proposition}[Endpoint contours in a square]
\label{q4pocket:square-cle}
In a homogeneous square, the entire signed contour hierarchy
converges to nested $\mathrm{CLE}_4$, with independent fair
signs at each generation.  The identification includes each
individual contour as a simple-loop element, its parent, and
its traversal order up to the loop convention.
\end{proposition}

\begin{proof}
The first-level set $A$ in Lemma~\ref{q4pocket:local-set}
has harmonic values $\pm2\lambda_0$.  Ordinary two-valued-set
uniqueness identifies its coupling with the field as the
labelled non-nested $\mathrm{CLE}_4$ gasket.

We must still identify the discrete contour elements, since
one such element could a priori have several winding
components.  Every positive contour has an interior-status
disk, by the two-status property and the exclusion of
boundary-only traversal.  Thus it winds at least one hole
of the identified gasket.  Select two distinct holes by
the first points they contain in a fixed countable dense
set.  Conditional on the gasket $A$, the two unit offsets
in our construction satisfy
\begin{equation}
 \mathbb E[v(\Omega_1)v(\Omega_2)\mid A]
 =\mathbb P\{\Omega_1\text{ and }\Omega_2
        \text{ belong to one contour}\mid A\}.
 \label{q4pocket:label-grouping}
\end{equation}
Indeed the signs were independent and fair given the entire
unlabelled contour geometry; different contours contribute
zero covariance, whereas the same contour contributes one.
In the unique two-valued-set coupling, different
$\mathrm{CLE}_4$ holes have independent fair labels, so the
left side of \eqref{q4pocket:label-grouping} is zero.
Countability excludes all grouping of distinct holes.

Consequently each positive contour winds exactly one simple
CLE hole.  Its full interior trace equals that hole's Jordan
boundary: the two-status property puts every trace point
on the boundary of its winding region, and winding stability
puts that boundary in the trace.  Additional boundary pieces
are excluded as in Proposition~\ref{q4pocket:regular-chordal}.
The limiting winding is one.  Any extra covering or
backtracking repeats a nondegenerate subarc and violates the
dimension-zero bound.  Thus the original contour is exactly
the corresponding simple-loop element, not just a set with
the same filled interior.

Proceed conditionally through the hierarchy.  Given the
geometry and labels through an identified level, the
residual fields in its holes are independent zero-boundary
GFFs by Proposition~\ref{q4pocket:conditional-field}.  The
next-level decomposition, with the parent constant subtracted,
makes its missing set in each hole a thin local set with
harmonic values $\pm2\lambda_0$.  Child winding components
lie in the parent by discrete winding stability at points
with clearance.  They are simply connected and cover almost
every point, by the arguments already given.  Every positive
child has an interior-status disk and is supported on the
closure of its winding regions by the two-status property.
Thus the same topology and thinness verification applies
inside each parent hole.

Conditional uniqueness, including measurability from the
respective residual field, gives independent fresh CLEs
in the parent holes.  Equation~\eqref{q4pocket:label-grouping}
and the trace argument again identify every individual child
and its sign.  Its parent can be read at any interior disk
point off the traces.  Induction identifies all finite
generations.  Finally every positive-diameter contour and
all its ancestors occur in the retained finite-cutoff
hierarchies, whose counts are tight.  Sending the cutoff
to zero identifies the full collection in the stated
matching topology.
\end{proof}

\begin{corollary}[Endpoint whole-plane ensemble]
\label{q4pocket:whole-plane}
The infinite-volume free critical FK loop ensemble at $q=4$
converges to the nested whole-plane $\mathrm{CLE}_4$ ensemble,
in the spherical matching metric for all loops.
\end{corollary}

\begin{proof}
Use the expanding-square coupling of Section~\ref{sec:loops}.
Its endpoint ingredients are bulk estimates: monotone
boundary coupling and exploration of actual boundary
connections couple the infinite-volume free switches with
the square switches deep inside a large square, except
when a boundary connection reaches the inner region.
Actual dual circuits on successive annuli block that
event, uniformly even under the maximal boundary law.
Actual-color annular blockers likewise make a loop from
a fixed disk reaching the coupling boundary unlikely.
The finite-square limit is given by
Proposition~\ref{q4pocket:square-cle}.

Nested simple CLEs in domains of conformal radius tending
to infinity converge locally to nested whole-plane CLE;
their spherical whole-plane law is M\"obius invariant
\citep[Theorem~1, Proposition~4, and Section~3.2]
{KemppainenWerner2016}.  Its range
$8/3<\kappa\leq4$ includes the present endpoint, and
centered expanding squares meet its conformal-radius
condition.

To pass from the local comparison to the spherical metric,
first fix a spherical diameter cutoff.  Loops lying
entirely outside a sufficiently large fixed disk are below
that cutoff.  With high probability, every other loop
above the cutoff remains inside the switch-coupling region,
by the preceding escape bound.  The number of such loops
is tight, and the finite-square matching identifies them
bijectively with their traversal order.  Take the mesh
limit, then enlarge the square and the fixed disk, and
finally send the diameter cutoff to zero.  This is the
required global matching statement, including all nesting
levels.
\end{proof}

\subsection{The prescribed marked Jordan approximations}

It remains to transfer Proposition~\ref{q4pocket:regular-chordal}
to the precise wired/free edge convention of the theorem.
This transfer uses actual color walls, so that a boundary
identification is never mistaken for a path in the lattice.

\begin{proposition}[Endpoint marked Jordan convergence]
\label{q4pocket:jordan-chordal}
For every marked uniform Jordan approximation specified in
the theorem, its critical $q=4$ distinguished curve converges
to chordal $\mathrm{SLE}_4$ in the limiting marked domain.
\end{proposition}

\begin{proof}
Use the two-change disk drawing from
Lemma~\ref{moving:drawing}: place tiles on every polygon-graph
edge, give boundary edges distinct exterior tabs, and put
local caps along the wired primal and complementary dual
arcs.  Flattening the tabs and separating the drawing in
a thin collar changes the ordered curve by only
$O(\delta_n)$ in uniform distance.  Boundary edges on the
wired primal arc may be fixed open, while edges on the
free arc retain their sampled law.

The primal bank of the distinguished path, closed along
the primal wall, is a walk of actual open edges in the
wired cluster.  On the other side, the actual internal
dual bank, its exits to exterior copies across the free
arc, and the local complementary-wire links give the
dual closing walk.  All artificial dual links and copies
stay within $O(\delta_n)$ of that free arc.  No dual
exit through the primal wired arc is used.  At fixed
bulk points away from limiting traces these walks
therefore have stable complementary winding statuses
$1$ and $0$.  Lemma~\ref{q4pocket:geometry} gives tightness,
the two-status property, and the restrictions on multiple
visits and boundary traversal for these drawings.

Construct a regular comparison polygon $V^+$ close to the
limiting marked domain $V$ with the following two strict
separations: the original primal wired arc, including
its endpoints, lies outside $\overline{V^+}$; and the
comparison dual wall lies outside $\overline V$.
All margins are fixed before taking the mesh limit.
One construction uses plane Schoenflies coordinates for
$\partial V$.  Move the primal portion slightly inward,
extending it a little beyond both marks, then rise
outside the unit circle near the two ends and put
both comparison transitions there.  Keep the whole
comparison dual portion outside the unit circle.
Approximate the resulting simple curve by a simple
rectilinear polygon, preserving all strict margins,
and place the transitions inside straight sides.
Its actual primal-open and dual-open wall approximations
are regular experiments.  Reverse the colors to construct
$V^-$: the original complementary arc is outside
$\overline{V^-}$, and the comparison primal wall is
outside $\overline V$.  Both comparisons approach the
same marked Jordan domain as the offsets vanish.

Couple the original law with that in $V^+$ so that
every original primal open edge in the common region
is open in the comparison.  To check the boundary
partition needed for this ordering, restrict the
original configuration to the common edges.  Any
connection departing from the comparison disk can
induce ties only at its actual primal wall vertices:
its dual wall and transitions lie outside the entire
original closure, while the original wired vertices
lie outside the comparison.  The comparison connects
all those primal wall vertices by fixed open wall
edges.  Its induced partition therefore dominates
the original one, and random-cluster monotonicity
gives the coupling.  The cut is along a primal
lattice wall, so no unaligned edge crossing is
being treated as a wall vertex.

Every piece inside $V^+$ of the original primal
closing bank walk belongs to the comparison primal
wall cluster: follow its original connection to
the outer wire and take its first exit from $V^+$.
Suppose a test point has original primal winding
status but comparison dual status, with positive
clearance.  On the comparison dual side it can be
joined to the comparison dual wall without crossing
the comparison primal wall cluster, and hence
without crossing the original closing walk inside
$V^+$.  From that wall, which is outside the entire
original polygon, it can escape to infinity through
the original exterior.  This contradicts the
original winding.  Thus original primal status
implies comparison primal status.

For $V^-$ apply exactly the dual ordering.  Original
dual identifications and artificial bank links lie
beside the complementary arc, outside $V^-$.  In
the common region the dual walk uses only internal
dual steps, and departing connections induce ties
only at the aligned comparison dual wall vertices.
An originally wired primal edge admits no dual
crossing step.  Original dual status therefore
implies comparison dual status.  The free-bank arm
bound removes near-hits and tile ambiguities at
any prescribed finite tuple of bulk points.

Let $I_n(z)$ be the original primal-side indicator,
and let $I_n^-(z),I_n^+(z)$ denote those in the
two comparisons.  The separate monotone couplings
give, for every fixed finite tuple $\boldsymbol z$,
the stochastic bounds
\begin{equation}
 (I_n^-(z))_{z\in\boldsymbol z}
 \ \preceq\
 (I_n(z))_{z\in\boldsymbol z}
 \ \preceq\
 (I_n^+(z))_{z\in\boldsymbol z},
 \label{q4pocket:side-sandwich}
\end{equation}
up to errors vanishing with the mesh.  At fixed
comparison polygons the limiting laws are the
regular chordal $\mathrm{SLE}_4$ laws by
Proposition~\ref{q4pocket:regular-chordal}.  Now let
the offsets vanish.  Conformal maps and their
boundary extensions converge under marked uniform
Jordan convergence, using uniform local connectivity
of these boundary parametrizations.  The two
comparison SLE laws therefore have the same limit.
Deterministic interior points avoid SLE, so their
side-status arrays converge as well.  Equation
\eqref{q4pocket:side-sandwich} identifies every
finite-dimensional status law of an original
subsequential limit.

On a fixed countable dense set, the statuses
determine the interior trace: an interior trace
point has both statuses in every neighborhood,
whereas a point off the trace has a small disk
of constant status.  This follows from
Lemma~\ref{q4pocket:geometry}.  Thus the limiting
interior trace is the SLE trace.  No boundary
traversal and the dimension-zero multiple-visit
bound identify the complete ordered curve exactly
as in Proposition~\ref{q4pocket:regular-chordal}.
Every tight subsequence has the same law, proving
the full-sequence assertion.
\end{proof}
